% =====================================================================
%  VARIABLES EDITABLES POR EL USUARIO
%  ---------------------------------------------------------------------
%  Rellena los valores siguientes con los datos de tu trabajo.
%  Edita solo el texto entre llaves { }; NO cambies el nombre de los
%  comandos, ya que los usan la portada y los metadatos del PDF.
% =====================================================================

% --- Información académica --------------------------------------------
% Nombre de tu grado / titulación.
\newcommand{\degree}{Grado en Ingeniería Informática}
% Curso académico (p. ej. 2024-2025).
\newcommand{\academicYear}{2024-2025}
% Tipo de trabajo; aparece en la portada y en los metadatos del PDF.
\newcommand{\thesisType}{Trabajo de Fin de Grado}

% --- Información del trabajo ------------------------------------------
% Título del trabajo.
\newcommand{\thesisTitle}{MI TÍTULO}
% Nombre completo del autor o autora.
\newcommand{\thesisAuthor}{AUTOR}
% Nombre completo del tutor o tutora.
\newcommand{\thesisAdvisor}{TUTOR}
% Lugar y fecha de la defensa (p. ej. Leganés, junio de 2025).
\newcommand{\placeAndDate}{LUGAR, MES AÑO}

% --- Institución -----------------------------------------------------
% Nombre de la universidad (se usa en los metadatos del PDF).
\newcommand{\university}{Universidad Carlos III de Madrid}

% --- Palabras clave --------------------------------------------------
% Aparecen en el Abstract / Resumen y se incrustan en los metadatos del
% PDF. Separa los términos con punto y coma.
\newcommand{\keywordsEn}{Keyword 1; Keyword 2; Keyword 3}
\newcommand{\keywordsEs}{Palabra clave 1; Palabra clave 2; Palabra clave 3}

% --- Configuración ---------------------------------------------------
% Licencia aplicada al trabajo. Valores admitidos: cc / reserved
\newcommand{\licenseType}{cc}
% Estilo de la bibliografía. Valores admitidos: apa / ieee
\newcommand{\bibStyle}{ieee}
% Mostrar el capítulo de demostración (chapters/00_examples) con ejemplos de la
% plantilla. Ponlo a false en la versión final (es solo contenido de ejemplo):
% así se oculta sin tener que comentar nada en tfg_main.tex.
\newcommand{\showDemoChapter}{true}
% (El interruptor PDF/A \pdfaOutput está al principio de tfg_main.tex, porque
%  \DocumentMetadata debe ejecutarse antes que todo lo demás.)
% (Los metadatos de la plantilla "Template.*" están en tfg_template_info.tex; los
%  mantiene quien actualiza la plantilla, no hace falta tocarlos para tu TFG.)


% =====================================================================
%  DECLARACIÓN DE USO DE IAG (anexo obligatorio)
%  ---------------------------------------------------------------------
%  La compone others/ai_declaration.tex. Los valores true/false
%  seleccionan y resaltan la opción correspondiente del formulario.
% =====================================================================

% Mostrar la ayuda, los ejemplos y la guía de tipos de uso (true) para una
% versión de trabajo, u ocultar todo lo no esencial (ayudas, comentarios,
% secciones no usadas) para una versión final lista para entregar (false).
% Ponlo a false antes de entregar.
\newcommand{\aiShowExplanations}{true}

% ¿Has usado IAG en tu TFG? (true/false)
% Si es false, se omiten las partes 1-3 y se marca "NO".
\newcommand{\aiUsed}{true}

% --- Parte 1: comportamiento legal, ético y responsable --------------
% Para cada pregunta: true = opción SÍ, false = opción NO.
\newcommand{\aiConfidentialData}{false}    % ¿datos de carácter confidencial?
\newcommand{\aiCopyrightedMaterial}{false} % ¿materiales con derechos de autoría?
\newcommand{\aiPersonalData}{false}        % ¿datos de carácter personal?
\newcommand{\aiTermsRespected}{true}       % ¿respetados términos de uso y ética?

% --- Parte 2: uso técnico --------------------------------------------
% Nombre(s) y versión(es) del sistema/herramienta de IAG usado.
\newcommand{\aiSystemName}{ChatGPT 4o}
% Una declaración de ejemplo (copia el bloque en others/ai_declaration.tex
% para declarar más usos).
\newcommand{\aiDeclPurpose}{buscar información}
\newcommand{\aiDeclPrompt}{Escribe aquí el prompt que utilizaste}
\newcommand{\aiDeclInteraction}{describe qué hiciste con la respuesta de la IA}

% --- Parte 3: reflexión sobre utilidad -------------------------------
\newcommand{\aiReflection}{Escribe aquí tu valoración personal (formato libre): las fortalezas y debilidades que has identificado al usar herramientas de IAG, y si te han servido en el aprendizaje, el desarrollo o las conclusiones.}


% =====================================================================
%  COLORES (avanzado)
%  ---------------------------------------------------------------------
%  Valores de color usados en toda la plantilla. La mayoría no necesitará
%  cambiarlos. Mantén el azul corporativo de la UC3M salvo que tu centro
%  permita otra cosa. Usa "R,G,B" (0-255) para RGB y hex de 6 dígitos para HTML.
%  Para AÑADIR un color nuevo tocas DOS ficheros: define aquí su valor y añade
%  el correspondiente \definecolor{<nombre>}{<modelo>}{\<tuvalor>} en tfg_uc3m.sty.
% =====================================================================
\newcommand{\colorBrand}{0,0,102}       % azul corporativo UC3M (portada, títulos) [RGB]
\newcommand{\colorHighlight}{FFFF00}    % resaltado de texto [HTML]
\newcommand{\colorLink}{0000EE}         % hiperenlaces, solo en modo PDF/A [HTML]
% Anexo de declaración de IA
\newcommand{\colorAiSelected}{FFE066}   % opción seleccionada [HTML]
\newcommand{\colorAiFill}{D7E9FF}       % valores que rellenas [HTML]
\newcommand{\colorAiHint}{1A7F5A}       % texto de ayuda [HTML]
% Listados de código. Los grises van de 0 = negro a 1 = blanco; el resto, HTML.
\newcommand{\colorCodeBg}{0.97}         % fondo (gris)
\newcommand{\colorCodeComment}{0.45}    % comentarios (gris)
\newcommand{\colorCodeNumbers}{0.55}    % números de línea (gris)
\newcommand{\colorCodeKeyword}{0033CC}  % palabras clave [HTML]
\newcommand{\colorCodeString}{008000}   % cadenas [HTML]
